
\section*{Foundations of Mixture of Experts}
\addcontentsline{toc}{section}{Foundations of Mixture of Experts}

\subsection*{The original MoE formulation}

The Mixture of Experts architecture was introduced by \cite{jacobs1991adaptive} as a modular neural network comprising multiple expert networks and a gating network that learns to partition the input space. Given an input $\mathbf{x}$, the MoE output is computed as:

\begin{equation}
    y = \sum_{i=1}^{N} g_i(\mathbf{x}) \cdot E_i(\mathbf{x})
\end{equation}

\noindent where $N$ is the number of experts, $E_i(\mathbf{x})$ denotes the $i$-th expert network, and $g_i(\mathbf{x})$ represents the gating weight assigned by the gating network $G$, typically implemented as a softmax over linear projections:

\begin{equation}
    g_i(\mathbf{x}) = \frac{\exp(\mathbf{w}_i^\top \mathbf{x})}{\sum_{j=1}^{N} \exp(\mathbf{w}_j^\top \mathbf{x})}
\end{equation}

This formulation enables the model to learn a soft partitioning of the input space, where different experts specialize in different regions or data patterns. The key theoretical advantage is that each expert needs to model only a subset of the data distribution, reducing the complexity of the learning problem for individual components while the ensemble maintains global expressiveness \cite{yuksel2012twenty}.

\subsection*{Sparse MoE and conditional computation}

The modern revival of MoE was catalyzed by the seminal work of \cite{shazeer2017outrageously}, who introduced the sparsely-gated MoE layer that activates only the top-$k$ experts for each input token. This sparse routing mechanism fundamentally changed the scalability equation: a model with $N$ experts and top-$k$ routing ($k \ll N$) can increase total parameter count by $N$-fold while increasing per-token computation by only $k/N$. The top-$k$ gating function is defined as:

\begin{equation}
    y = \sum_{i \in \text{TopK}(G(\mathbf{x}), k)} g_i(\mathbf{x}) \cdot E_i(\mathbf{x})
\end{equation}

\noindent where $\text{TopK}(\cdot, k)$ selects the $k$ experts with the highest gating scores. This approach enabled the training of models with trillions of parameters, such as the Switch Transformer \cite{fedus2022switch}, while maintaining inference costs comparable to much smaller dense models.

\subsection*{Load balancing and expert utilization}

A critical challenge in MoE training is ensuring balanced utilization of all experts. Without explicit regularization, the gating network tends to converge to a state where a small number of experts receive the majority of inputs (``rich get richer'' dynamics), leading to expert collapse where most experts remain undertrained. To address this, auxiliary load-balancing losses have been proposed:

\begin{equation}
    \mathcal{L}_{\text{balance}} = N \cdot \sum_{i=1}^{N} f_i \cdot P_i
\end{equation}

\noindent where $f_i$ is the fraction of tokens dispatched to expert $i$ and $P_i$ is the average gating probability for expert $i$ \cite{fedus2022switch}. Recent innovations include expert-level dropout \cite{zoph2022st}, capacity factors that limit the maximum number of tokens per expert, and hash-based routing that guarantees uniform assignment \cite{roller2021hash}.

\subsection*{MoE in the transformer era}

The integration of MoE into transformer architectures \cite{vaswani2017attention} has been the primary driver of recent advances. In transformer-based MoE models, the feed-forward network (FFN) layers in selected transformer blocks are replaced with MoE layers, where each ``expert'' is itself an FFN. This design preserves the attention mechanism for contextual token mixing while introducing conditional computation in the position-wise transformation. Key architectural variants include:

\textbf{Switch Transformer} \cite{fedus2022switch}: simplifies MoE routing to top-1 (each token sent to a single expert), dramatically simplifying the design while maintaining performance.

\textbf{GLaM} \cite{du2022glam}: demonstrates that a 1.2 trillion parameter MoE model can outperform GPT-3 \cite{brown2020gpt3} while using only one-third of the training energy.

\textbf{Mixtral 8$\times$7B} \cite{jiang2024mixtral}: employs 8 experts per layer with top-2 routing, achieving performance comparable to dense models with 4$\times$ more parameters.

\textbf{DeepSeek-V2/V3} \cite{deepseek2024, zhang2024deepseek_v3}: introduces fine-grained expert segmentation with shared experts, pushing expert specialization to its limits.

\textbf{LLaMA-4} \cite{llama4_2025}: integrates MoE into Meta's open-source LLM family with up to 128 experts per layer.

These advances have established MoE as the de facto architecture for scaling large language models, and the resulting techniques are now being adapted for healthcare-specific applications across multiple modalities and tasks.
